\section{The two-scale secondary saturation domain}\label{corner:section}

We determine the coefficient region reached when three roots of the
principal quartic collapse to zero. A division of the slope gives the finer
necessary scales; the joint return quotient and root charts then identify
the region and realize its interior.
Throughout this section we use the setting and conventions of
Section~\ref{sup:setting}.

\subsection{The quartic corner and its cubic limit}

The principal-scale coordinates
of~\cite[Section~\ref*{main-subsec:four-root-limit}]{main2026} are
\[
 A_4=\frac{C^3a}{m^4},\qquad S_4=\frac{Ch}{m^2},\qquad
 T_4=\frac{3BC^2\delta}{2m^3},\qquad x=\frac{C\sqrt q}{m},
 \qquad h=-b.
\]
Write $e_j$ for the elementary symmetric polynomial of degree $j$, and set
\[
 \Phi_{A,S,T}(x)=A-Sx^2+x^3-x^4+Tx,\qquad
 \mathcal K_4=\{(-e_4,e_2,e_3):r_i\geq0,\ \textstyle\sum_{i=1}^4r_i=1\}.
\]
Thus $\Phi_{A_4,S_4,T_4}=-\prod_{i=1}^4(x-r_i)$ for a point in
$\mathcal K_4$. We abbreviate this polynomial by $\Phi$.
The corner $S_4\to0$ of $\mathcal K_4$ explains the next scale.
Order the four roots and write $u=r_1+r_2+r_3$, $v=r_1r_2+r_1r_3+r_2r_3$,
$w=r_1r_2r_3$. Since $r_4=1-u$,
\[
 S_4=(1-u)u+v,\qquad T_4=(1-u)v+w,\qquad -A_4=(1-u)w.
\]
Along any sequence with $S_4\to0^+$, the largest root tends to one and
$u/S_4\to1$. After passing to a subsequence, $r_i/S_4\to z_i\ge0$ for
$i=1,2,3$, with $\sum_{i=1}^3z_i=1$. Along this subsequence,
$(A_4/S_4^3,T_4/S_4^2)\to(-e_3(z),e_2(z))$.
Conversely, taking small roots $t z_i$ with $\sum z_i=1$
and large root $1-t$ realizes every such limit; here
$S_4=t-t^2(1-e_2(z))\sim t$.

Set $\sigma_4=S_4$ and write $A_4=\sigma_4^3\alpha$,
$T_4=\sigma_4^2T_3$, and $x=\sigma_4z$. Then
\[
 \Phi(\sigma_4z)/\sigma_4^3
 =\alpha-z^2+z^3+T_3 z-\sigma_4z^4.
\]
The polynomial identity identifies a candidate cubic. Transferring it to
the return map when $h=-b=o(m^2)$ requires a division before either scale
vanishes: an additive error $o(m^2)$ in the slope would not control $h$.
The necessary estimates use sequences whose original parameters, including
all remaining parameters, tend to the reference of Section~\ref{sup:setting}.

\subsection{Slope division and necessary bounds}

The finer scales are forced by the slope at the first zero of $P''$ in the
outer region and by the extrema of the displacement. We begin with a division
of the slope on the separation graph.

For a return with six fixed points, let $a_\delta=\Gamma_1^{-1}(\delta)$ be the
second threshold. By~\cite[Lemma~\ref*{main-sat:saturation}]{main2026} each
open curvature block contains one zero $t_j$ of $D_q$, with extremal signs
$+,-,+,-,+$, and
\begin{equation}\label{corner:saturation-slopes}
 D_q(0)<0,\quad D_q(a_\delta)>0,\quad D_q(\zeta_1)<0,\quad D_q(\zeta_2)>0.
\end{equation}
Here \(\zeta_1<\zeta_2\) are the two zeros of \(P''\) in the outer region; in
particular
\(m,\delta,h=-b\) are positive.

The separation graph of~\cite[Section~3]{main2026} is
\begin{equation}\label{corner:curvature-division}
 \begin{aligned}
 \delta=\Delta(r,m,p)&=r^2\{mH_\Delta(r,m,p)+rB_\Delta(r,p)\},\\
 H_\Delta(0,0,p)&=\frac2{A(p)},\qquad
 B_\Delta(0,p)=\frac{16c_0(p)}{3k(p)A(p)}.
 \end{aligned}
\end{equation}
Here $B_\Delta,H_\Delta$ are smooth profiles.

\begin{lemma}[Slope on the separation graph]
\label{corner:graph-slope}
After a common reduction there are jointly smooth profiles $H_D,B_D$
such that
\begin{equation}\label{corner:exact-graph-slope}
 \left.D_q(r^2,p,m,\delta)\right|_{\delta=\Delta(r,m,p)}
 =b+r\{mH_D(r,m,p)+rB_D(r,p)\},
\end{equation}
with
\begin{equation}\label{corner:graph-slope-coefficients}
 H_D(0,0,p)=3k(p),\qquad B_D(0,p)=6c_0(p).
\end{equation}
\end{lemma}

\begin{proof}
Use the signed-root profiles
of~\cite[Section~3]{main2026}.
With $\delta=r^2v$, the expression
\[
 D_q(r^2)=R'(Q)\ell_2(r\sqrt{T_1(r)^2-v})\ell_1(r)-1
\]
extends smoothly in $(r,m,v,p)$ on a common box.  At $r=0$ it equals
$b$ for all small $m,v$.  Hadamard division in $r$ therefore gives a
smooth function $\mathcal T$ such that $D_q-b=r\mathcal T$.
The first passage derivatives give
\[
 \ell_\nu(u)=1+\tfrac32 A_\nu u+O(u^2).
\]
Consequently, at the root boundary,
\[
 \mathcal T(0,m,v,p)
 =\tfrac32 k(p)
   \{A_1(p,m,0)+A_2(p,m,0)\sqrt{1-v}\}.
\]
The aligned strength identity yields the linear terms
$\tfrac32km+\tfrac34kAv$.  On $m=v=0$, finite cancellation gives
$D_q-b=2c_0r^2+O(r^3)$.  Joint Taylor division thus gives
\begin{equation}\label{corner:fixed-control-slope-expansion}
 \mathcal T
 =2c_0r+\tfrac32km+\tfrac34kAv
       +O\bigl((|r|+|m|+|v|)^2\bigr),
\end{equation}
with a uniform smooth remainder.

Now substitute the graph
$v=V(r,m,p)=mH_\Delta+rB_\Delta$ in $\mathcal T$, and call the resulting
smooth function $\mathcal T_\Delta(r,m,p)$. Formula
\eqref{corner:fixed-control-slope-expansion} gives
\[
 \mathcal T_\Delta(0,0,p)=0,\qquad
 \partial_m\mathcal T_\Delta(0,0,p)=3k,\qquad
 \partial_r\mathcal T_\Delta(0,0,p)=6c_0.
\]
Define $B_D(r,p)=\mathcal T_\Delta(r,0,p)/r$ by smooth division and
\[
 H_D(r,m,p)=\int_0^1\partial_m\mathcal T_\Delta(r,sm,p)\,ds.
\]
Thus $\mathcal T_\Delta=rB_D+mH_D$, proving
\eqref{corner:exact-graph-slope}--\eqref{corner:graph-slope-coefficients}.
\end{proof}

For $\delta>0$, $q\geq0$, the global passage estimate
\eqref{sup:global} gives
\begin{align}
 |D(q)-a|\leq K\{&
 hq+q^2+mq^{3/2}+m^2q^2+q^{5/2}\notag\\
 &+\delta\sqrt{q+\delta}+\delta q^{3/2}\},
 \qquad h=-b>0.                         \label{corner:global-comparison}
\end{align}
Indeed $0\leq Q(q)\leq Kq$ uniformly.  Below $a_\delta$ this follows from
$Q=\Gamma_1(q)$; above $a_\delta$, positivity and a common bound
$\Gamma_2(y)\leq Ky$, with $K\geq1$, give
$\delta+\Gamma_2(\Gamma_1(q)-\delta)\leq K\Gamma_1(q)$.
Insert this bound in the closing expansion
\eqref{sup:closing} and apply
\eqref{sup:global}.

\begin{proposition}[Secondary necessary bounds]
\label{corner:secondary-bounds}
\label{corner:sharp-positive-offset}
\label{corner:positive-offset-bound}
There is a constant $K$ such that, for parameters sufficiently close to the
reference in this sign case, every return with six fixed points satisfies,
with $h=-b$,
\begin{equation}\label{corner:strong-bounds}
 0<h\leq Km^2,\qquad
 0<\delta\leq K\frac{h^2}{m},\qquad
 -K\frac{h^3}{m^2}\leq a\leq Kh^3.
\end{equation}
More precisely, with $k=1-h$ and the current first passage strength
$A_1(\lambda)$,
\begin{equation}\label{corner:finite-positive-offset-bound}
 a<\frac{4h^3}{27k^2A_1(\lambda)^2}+Kh^4.
\end{equation}
Along every six-fixed-point sequence tending to the reference, this gives
\begin{equation}\label{corner:sharp-positive-offset-limsup}
 \limsup\frac{a}{h^3}\leq\frac{4}{27B^2}.
\end{equation}
\end{proposition}

The leading term in~\eqref{corner:finite-positive-offset-bound} agrees with
Simpson and Meiss's fold value $4(1-\Xi)^3/(27\gamma^2)$ for the truncated
return~\cite[arXiv version, Eq.~(72)]{simpsonmeiss2008}, with $h=1-\Xi$ and
$\gamma=kA_1(\lambda)$. Here the bound retains $Kh^4$ and is necessary for
arbitrary six-cycle configurations; Proposition~\ref{corner:positive-offset-realization}
proves sharpness by sequences approaching the endpoint.

\begin{proof}
\cite[Proposition~\ref*{main-sat:weighted-confinement}]{main2026} gives
\(h\le Km^2\).
Put \(\rho_i=\sqrt{\zeta_i}\).

The root coordinate $\rho_1$ of the first curvature zero lies below the
graph's unique positive critical point, whose expansion is
\[
 \rho_1<r_{\rm c}
 =-\frac{k(p)m}{4c_0(p)}+O(m^2).
\]
Uniformly for $0<r\leq r_{\rm c}$, the slope profiles satisfy
\[
 mH_D+rB_D\geq c_*m,\qquad c_*>0.
\]
To check the margin, divide by $m$ and use
\eqref{corner:graph-slope-coefficients}: the limiting minimum on this
interval is
$3k+6c_0[-k/(4c_0)]=3k/2>0$.
The negative slope at $\zeta_1$ therefore yields
\[
 \rho_1\leq K\frac hm,\qquad
 \zeta_1\leq K\frac{h^2}{m^2}.
\]
Using \eqref{corner:curvature-division} at $\rho_1$, where
$mH_\Delta+\rho_1B_\Delta=O(m)$, proves
$\delta\leq Kh^2/m$.

For the upper offset bound, the negative minimum $t_2=r_2^2$ lies
where only the first passage is active, so that $Q=\Gamma_1(q)$. Uniform
passage and closing expansions at fixed parameters give
\[
 D(r^2)=a-hr^2+(1-h)A_1r^3+O(r^4),\qquad
 D_q(r^2)=-h+\tfrac32(1-h)A_1r+O(r^2).
\]
The coefficient $\tfrac32(1-h)A_1$ has a uniform positive lower bound.
After reducing the state interval, the $O(r_2^2)$ term is absorbed in
that linear term. Thus $D_q(t_2)=0$ first gives $r_2=O(h)$ and then
\[
 r_2=\frac{2h}{3kA_1(\lambda)}+O(h^2).
\]
The value expansion therefore yields
\begin{equation}\label{corner:first-minimum-sharp}
 D(t_2)=a-\frac{4h^3}{27k^2A_1(\lambda)^2}+O(h^4).
\end{equation}
Its strict negativity proves \eqref{corner:finite-positive-offset-bound}
and the coarse upper bound in \eqref{corner:strong-bounds}.
Since $k\to1$ and $A_1(\lambda)\to B$ along a converging sequence,
it also proves \eqref{corner:sharp-positive-offset-limsup}.
Thus $a\leq4h^3/(27B^2)+o(h^3)$.

For the lower bound use the third extremum $t_3\in(a_\delta,\zeta_1)$.
It has $D(t_3)>0$ and $t_3\leq Kh^2/m^2$.  Substitute this and
$\delta\leq Kh^2/m$ in \eqref{corner:global-comparison}.
The terms $ht_3$ and $mt_3^{3/2}$ are $O(h^3/m^2)$.
The term $t_3^2$ is bounded by that size times $h/m^2$, which is
bounded by the first estimate.  The higher aligned terms are smaller.
Finally,
\[
 \delta\sqrt{t_3+\delta}
 \leq K\left(\frac{h^3}{m^2}+\frac{h^3}{m^{3/2}}\right).
\]
The remaining separation term is smaller as well. Thus
$|D(t_3)-a|\leq Kh^3/m^2$, proving the lower offset bound.

\end{proof}


\subsection{The joint secondary quotient}

In the standard unfolding introduce
\begin{equation}\label{corner:weights}
 m=\mu,\qquad b=-\mu^2\sigma,\qquad
 \delta=\mu^3\sigma^2d,\qquad a=\mu^4\sigma^3\alpha,
 \qquad \mu,\sigma>0.
\end{equation}
These weights give $\sigma_4=S_4=C\sigma$ and $T_3=3Bd/2$.
The weight map has determinant $-\mu^9\sigma^5$ with rows
$(m,b,a,\delta)$ and columns $(\mu,\sigma,\alpha,d)$.
Here $h=\mu^2\sigma$ and the secondary positive state coordinate is
\[
 q=\mu^2\sigma^2z^2.
\]
The threshold layer has relative size $\delta/q=\mu d/z^2$, so on fixed
positive $z$-intervals both passages are active for small $\mu$, even when
$\sigma$ tends to zero independently.

\begin{definition}[Secondary saturation domain]
\label{def:actual-secondary-domain}
Use the neighbourhood $\mathcal V$, the standard unfolding $\Lambda$ and the
count $N_\lambda(I)$ of Section~\ref{sup:setting}.
Let $w_{\mu,\sigma}(\alpha,d)=(a,b,m,\delta)$ be given by
\eqref{corner:weights}. The secondary saturation domain is
\[
 \mathscr S_{\mu,\sigma}
 =\{(\alpha,d)\in\R^2:w_{\mu,\sigma}(\alpha,d)\in\mathcal V,
       N_{\Lambda(w_{\mu,\sigma}(\alpha,d))}(I)=6\}.
\]
\end{definition}

Define the compact set
\begin{equation}\label{corner:CB}
 \mathcal C_B=
 \left\{\left(-r_1r_2r_3,\,
       \frac{2(r_1r_2+r_1r_3+r_2r_3)}{3B}\right):
       r_i\geq0,\quad r_1+r_2+r_3=1\right\}.
\end{equation}

We use Euclidean distance for the Hausdorff distance $d_H$ between compact
coefficient sets. The quotient below connects the cubic region to the return
map, with no relation between $\mu$ and $\sigma$. Further charts and
estimates of critical points then account for the remaining roots and the
boundary.

\begin{lemma}[Joint secondary quotient]
\label{corner:joint-quotient}
On every compact positive $z$-interval and bounded set of rescaled
parameters, the function
\[
 \mathcal F_{\mu,\sigma}(z;\alpha,d)
 =\frac{D(\mu^2\sigma^2z^2)}{\mu^4\sigma^3}
\]
extends jointly smoothly to signed $(\mu,\sigma)$ near $(0,0)$. There is a
jointly smooth profile $E$ such that
\begin{equation}\label{corner:joint-limit}
 \mathcal F_{\mu,\sigma}
 =\alpha-z^2+z^3+\tfrac32 Bdz-C\sigma z^4
       +\mu E(\mu,\sigma,z,\alpha,d).
\end{equation}
Every prescribed finite collection of mixed derivatives is bounded on
closed reductions. The extension equals the rescaled displacement for
$\mu,\sigma>0$.
\end{lemma}

\begin{proof}
The root-profile division
\eqref{sup:root-division}, before any weights are imposed, gives
\begin{equation}\label{corner:exact-Q-division}
 Q-q=mr^3+\tfrac34m^2r^4+r^5L_0+\delta rH(r,m,v,p),
 \qquad v=\delta/r^2,\qquad H(0,0,0,p_0)=\tfrac32B.
\end{equation}
Set $r=\mu\sigma z$ and $v=\mu d/z^2$.  Define
\begin{align}
 V&=z^3+dzH+\mu\sigma^2z^5L_0
                    +\tfrac34\mu^2\sigma z^4,\notag\\
 Y&=z^2+\mu^2\sigma V.                  \label{corner:VY}
\end{align}
These are jointly smooth for positive $z$ in a compact interval,
including signed auxiliary $\mu,\sigma$, and satisfy
$Q-q=\mu^4\sigma^3V$, $Q=\mu^2\sigma^2Y$ for positive scales.
The closing expansion gives
\begin{equation}\label{corner:exact-residual}
 \mathcal F_{\mu,\sigma}
 =\alpha+V-Y+c(\lambda)\sigma Y^2
       +\mu^2\sigma^3Y^3S(\mu^2\sigma^2Y,\lambda).
\end{equation}
At $\mu=0$, all four controls in
\eqref{corner:weights} vanish, for every small $\sigma$.
Because $\xi$ is fixed, $c=-C$ and $H=3B/2$ there.  The restriction of
\eqref{corner:exact-residual} to $\mu=0$ is therefore exactly
$\alpha-z^2+z^3+(3B/2)dz-C\sigma z^4$.
Taylor integral division in $\mu$ proves \eqref{corner:joint-limit}
and the mixed bounds.
\end{proof}

For a dependent family whose original parameters tend to the reference,
if $\mu_n,\sigma_n\to0^+$ and $(\alpha_n,d_n)$ stay bounded, the same
identity gives, on every compact positive $z$ interval,
\begin{equation}\label{corner:sequence-limit}
 \frac{D_n(\mu_n^2\sigma_n^2z^2)}
           {\mu_n^4\sigma_n^3}
 -\{\alpha_n-z^2+z^3+\tfrac32Bd_nz\}\longrightarrow0
\end{equation}
in every prescribed finite state derivative.  Indeed
\eqref{corner:exact-Q-division} and \eqref{corner:exact-residual} still
hold along the sequence, with its original remaining parameters retained.
Their profiles converge to their reference values, and all their
$z$ derivatives are controlled by the same smooth root expressions.

\subsection{Charts for the remaining roots}

The following charts supply the remaining roots both in the secondary
interior and in the positive-offset configurations of
Section~\ref{sup:positive}.

\begin{lemma}[Charts for the remaining roots]
\label{corner:realization-charts}
In the standard unfolding, for bounded $\alpha,d$ with $d>0$, the following
limits hold in $C^1$ on the indicated compact
state intervals as $\mu,\sigma\to0^+$ independently:
\begin{align}
 \frac{D(\mu^2\sigma^2x)}{\mu^4\sigma^3}
     &\longrightarrow\alpha-x,
       &&x\Subset(-\infty,0),                         \label{corner:negative-small-chart}\\
 \frac{D(-hX/C)}{h^2/C}
     &\longrightarrow X-X^2,
       &&X\Subset(0,\infty),                          \label{corner:negative-large-chart}\\
 \frac{D(\mu^2X)}{\mu^4}
     &\longrightarrow X^{3/2}-CX^2,
       &&X\Subset(0,\infty).                          \label{corner:positive-large-chart}
\end{align}
Convergence is uniform when the rescaled parameters range over a compact
set.
\end{lemma}

\begin{proof}
For $q<0$ both passages are exactly inactive because $\delta>0$.
Thus \eqref{sup:closing} gives
$D(q)=a+bq+c(\lambda)q^2+q^3S(q,\lambda)$.
Substituting $q=\mu^2\sigma^2x$ and $q=-hX/C$, using
\eqref{corner:weights}, $h=\mu^2\sigma$ and $c(\lambda)\to-C$,
proves \eqref{corner:negative-small-chart} and
\eqref{corner:negative-large-chart}.

For the positive outer chart set $r=\mu\sqrt X$ and
$v=\mu\sigma^2d/X$ in the root division
\eqref{corner:exact-Q-division}. Here $v\to0$ uniformly on the prescribed
positive compact intervals, so both passages are active in this chart. The division gives
$(Q-q)/\mu^4\to X^{3/2}$ and $Q/\mu^2\to X$ in $C^1$.
In \eqref{sup:closing}, $a/\mu^4=\sigma^3\alpha$,
$b/\mu^2=-\sigma$, $c(\lambda)\to-C$, and the cubic closing remainder
divided by $\mu^4$ tends to zero. This proves
\eqref{corner:positive-large-chart}. Smooth profiles give the stated
$C^1$ uniformity on compact sets of states and rescaled parameters as
$\mu,\sigma\to0^+$ independently.
\end{proof}

\subsection{Critical points and the limiting cubic}

For the necessary limits we again allow dependent controls.
For $d\ge0$, write $T_3=3Bd/2$ and
\[ F_{\alpha,d}(z)=z^3-z^2+T_3 z+\alpha. \]

Necessity requires control at the left endpoint of the positive root
chart. We first place the third and fourth extrema on the secondary
scale. For a positive limit of $d$, separated passage slopes exclude
the threshold layer; when $d\to0$, the third maximum supplies the missing
value estimate.

\begin{lemma}[Limiting critical points]
\label{corner:limiting-critical-states}
Let a sequence in an admissible family tend to its reference,
with six distinct fixed points, $\mu=m\to0^+$ and
$\sigma=-b/m^2\to0^+$. Pass to a subsequence on which the bounded
normalized parameters satisfy $(\alpha,d)\to(\alpha_0,d_0)$.
If $d_0>0$, the third and fourth displacement extrema satisfy
\[
 \frac{\sqrt{t_3}}{\mu\sigma}\longrightarrow
 z_-=\frac{1-\sqrt{1-3T_{3,0}}}3,\qquad
 \frac{\sqrt{t_4}}{\mu\sigma}\longrightarrow
 z_+=\frac{1+\sqrt{1-3T_{3,0}}}3,
 \qquad T_{3,0}=\tfrac32Bd_0,
\]
where $0<T_{3,0}\leq1/3$.  Moreover,
\begin{equation}\label{corner:critical-value-inequalities}
 F_{\alpha_0,d_0}(z_-)\geq0,\qquad
 F_{\alpha_0,d_0}(z_+)\leq0,\qquad \alpha_0\leq0.
\end{equation}
If $d_0=0$, then $\alpha_0=0$.
\end{lemma}

\begin{proof}
Proposition~\ref{corner:secondary-bounds} gives
\[
 -K\leq\alpha\leq K\mu^2,\qquad 0<d\leq K.
\]
Write $\rho_1=\mu\sigma z_c$.  The same bounds control $z_c$.
The graph and slope identities yield
\begin{equation}\label{corner:normalized-curvature}
 d=z_c^2\{H_\Delta+\sigma z_cB_\Delta\}
       =\frac2Bz_c^2+o(1),\qquad
 \frac{D_q(\zeta_1)}h=-1+z_cH_D+\sigma z_c^2B_D
       =-1+3z_c+o(1).
\end{equation}
The strict negative slope gives $\limsup z_c\leq1/3$ and therefore
$d_0\leq2/(9B)$.

We already have $0<t_3<\zeta_1=O(\mu^2\sigma^2)$.
To obtain the same bound for $t_4$, first note that
\begin{equation}\label{corner:larger-curvature-radius}
 \rho_2/\mu\longrightarrow\frac3{8C}.
\end{equation}
In fact $\rho_2=O(\mu)$ by the curvature graph, it is beyond the
positive critical point $r_{\rm c}/\mu\to1/(4C)$ of the graph, and
$\delta/\mu^3=\sigma^2d\to0$.  Dividing
\eqref{corner:curvature-division} by $\mu^3$ shows that every limit
$w$ of $\rho_2/\mu$ satisfies
\[
 0=w^2\left(\frac2B-\frac{16C}{3B}w\right).
\]
The lower bound supplied by this critical point selects $w=3/(8C)$.
Choose a fixed $Z$ larger than the bound for $z_c$ and large enough
that $-1+3Z/2>0$.  By \eqref{corner:sequence-limit} and the chain rule in the state,
\[
 \frac{D_q(\mu^2\sigma^2Z^2)}h
 =-1+\tfrac32Z+\frac{3Bd}{4Z}+o(1)>0.
\]
For small scales this state lies in $(\zeta_1,\zeta_2)$, by
\eqref{corner:larger-curvature-radius}.  The derivative is strictly
increasing on that interval, so its unique zero $t_4$ lies before
the test state.  Thus both normalized critical states are bounded.

Suppose now that $d_0>0$. We exclude a zero of $D_q$ at
$q=\mu^2\sigma^2z^2$ with $z\to0$ where both passages are active.
Put $y=\Gamma_1(q)-\delta\geq0$. Along this sequence the parameters
and state lie in the common small domain, $A_1\ge B/2$, and $m,\delta>0$.
Since
\[
 \frac{q^{3/2}}{\delta}=\frac{\sigma z^3}{d}\longrightarrow0,
\]
the estimate of~\cite[(\ref*{main-sat:threshold-slope-exclusion})]{main2026} applies and gives
$D_q(q)\geq\kappa\delta/\sqrt q-K(h+q)$.
With the exact substitutions
\[
 \frac{\delta}{\sqrt q\,h}=\frac dz,\qquad
 \frac qh=\sigma z^2,
\]
we obtain, for a constant $c_*>0$,
\begin{equation}\label{corner:event-layer-slope-exclusion}
 \frac{D_q(q)}h\geq-K+c_*\frac dz-K\sigma z^2.
\end{equation}
The right side tends to $+\infty$, which contradicts a zero of $D_q$. The
estimate also holds as $y\to0$, at the second threshold.

The two normalized critical points now remain in a compact subset of
$(0,\infty)$.  Their limits are zeros of
$F'_{\alpha_0,d_0}=3z^2-2z+T_{3,0}$, by
\eqref{corner:sequence-limit}.  Their order on either side of $z_c$,
whose limit is $\sqrt{Bd_0/2}=\sqrt{T_{3,0}/3}$,
selects $z_-,z_+$, respectively; at $T_{3,0}=1/3$ the two
limits agree.  The signs $D(t_3)>0$, $D(t_4)<0$ give
\eqref{corner:critical-value-inequalities}.  The upper offset bound
already gives $\alpha_0\leq0$.

If $d_0=0$, \eqref{corner:normalized-curvature} instead gives
$z_c\to0$, so $z_M=\sqrt{t_3}/(\mu\sigma)\leq z_c\to0$.
At this state \eqref{corner:global-comparison}, divided by
$\mu^4\sigma^3$, gives the explicit bound
\begin{align*}
 \frac{|D(t_3)-a|}{\mu^4\sigma^3}\leq K\{&
 z_M^2+z_M^3+\sigma z_M^4+\mu^2\sigma z_M^4
       +\mu\sigma^2z_M^5\\
 &+dz_M+\sqrt\mu\,d^{3/2}
       +\mu^2\sigma^2d z_M^3\}.
\end{align*}
Every term tends to zero.  Positivity of $D(t_3)$ implies
$\alpha_0\geq0$, and $\alpha\leq K\mu^2$ implies
$\alpha_0\leq0$.  Therefore $\alpha_0=0$.
\end{proof}

\subsection{The complete corner domain}

The limits of the critical points give necessity, while the joint quotient
and the charts of Lemma~\ref{corner:realization-charts} give interior
realization.
The interior of $\mathcal C_B$ consists precisely of parameters for which
$F_{\alpha,d}$ has three distinct positive roots.  Equivalently,
\begin{equation}\label{corner:discriminant}
 \mathcal C_B=
 \left\{(\alpha,d):
 \begin{array}{l}
 \alpha\leq0,\quad d\geq0,\\
 T_3^2-4T_3^3+4\alpha-18T_3\alpha-27\alpha^2\geq0
 \end{array}\right\}.
\end{equation}

\begin{proposition}[Secondary corner domain]
\label{corner:complete-cluster-set}
\label{thm:saturation-geometry}
\label{cor:saturation-domain-limit}
In the fixed sign sector the following statements hold.
\begin{enumerate}
\item
Let parameters in an admissible family tend to the reference.  Suppose each return has six distinct fixed points and
\[
 m\longrightarrow0^+,\qquad
 \sigma=\frac{-b}{m^2}\longrightarrow0^+.
\]
Then the normalized parameters
\begin{equation}\label{corner:normalized-observations}
 \alpha=\frac{a}{m^4\sigma^3}
        =\frac{am^2}{h^3},
 \qquad
 d=\frac{\delta}{m^3\sigma^2}
   =\frac{m\delta}{h^2},\qquad h=-b,
\end{equation}
are bounded, and every cluster point belongs to $\mathcal C_B$.
In particular
\begin{equation}\label{corner:sharp-cluster-bounds}
 \liminf\alpha\geq-\frac1{27},\qquad
 \limsup\alpha\leq0,\qquad
 0\leq\liminf d\leq\limsup d\leq\frac2{9B}.
\end{equation}
If $a\geq0$ eventually, the last upper bound improves to
\begin{equation}\label{corner:nonnegative-offset-bound}
 \limsup d\leq\frac1{6B}.
\end{equation}

\item
Assume~\cite[(\ref*{main-eq:rank})]{main2026} and use the standard
unfolding. For every compact
$\mathcal K\Subset\operatorname{int}\mathcal C_B$, there exist
$\mu_0,\sigma_0>0$ such that, for every
\[
 0<\mu<\mu_0,\qquad 0<\sigma<\sigma_0,\qquad
 (\alpha,d)\in\mathcal K,
\]
the parameters \eqref{corner:weights} give exactly six simple fixed points
in $I$.  If $r_1<r_2<r_3$ are the positive roots of
$F_{\alpha,d}$, their states are, uniformly on $\mathcal K$,
\begin{align}
 q_{N,1}&=-\frac hC(1+o(1)),&
 q_{N,2}&=\mu^2\sigma^2(\alpha+o(1)),\notag\\
 q_{P,j}&=\mu^2\sigma^2(r_j^2+o(1)),\quad 1\leq j\leq3,&
 q_{P,4}&=\frac{\mu^2}{C^2}(1+o(1)).
                                      \label{corner:interior-roots}
\end{align}
Here $o(1)\to0$ as $\mu+\sigma\to0$, with the two positive scales
independent.

\item
Every point of $\mathcal C_B$, including every boundary point, is
the limit of normalized parameters along a sequence of returns with six
fixed points in the standard unfolding.

The unrestricted bounds in \eqref{corner:sharp-cluster-bounds} are sharp.
At the top endpoint $d_0=2/(9B)$, necessarily
$\alpha_0=-1/27$, and the limiting cubic is $(z-1/3)^3$.

For the standard unfolding, $\mathscr S_{\mu,\sigma}$ has nonempty compact
closure in one common rectangle for all sufficiently small independent
positive $\mu,\sigma$, and
\begin{equation}\label{eq:main-Hausdorff}
 d_H(\overline{\mathscr S_{\mu,\sigma}},\mathcal C_B)\longrightarrow0
 \quad(\mu+\sigma\to0).
\end{equation}
\end{enumerate}
\end{proposition}

Figure~\ref{fig:corner-limit-set} depicts $\mathcal C_B$.

\begin{figure}[htbp]
\centering
% Exact limiting set in coordinates (alpha,B d); not a finite-scale atlas.
\begin{tikzpicture}[x=1cm,y=1cm,>=Stealth,font=\small]
\fill[black!9] plot[domain=0:0.5,samples=181,variable=\t]
 ({-120*\t*\t*(1-2*\t)},{(40/3)*(2*\t-3*\t*\t)}) -- (0,0) -- cycle;
\draw[->] (-5.05,0)--(0.85,0) node[right]{$\alpha$};
\draw (0,-0.25)--(0,0);
\draw[->] (0,{10/3})--(0,5.05) node[above]{$B d$};
\draw[thick] plot[domain=0:0.5,samples=181,variable=\t]
 ({-120*\t*\t*(1-2*\t)},{(40/3)*(2*\t-3*\t*\t)});
\draw[thick,densely dashed] (0,0)--(0,{10/3});
\draw[black!50,densely dotted] ({-40/9},{40/9})--(0,{40/9});
\draw[black!50,densely dotted] ({-40/9},0)--({-40/9},{40/9});
\draw ({-40/9},0.06)--({-40/9},-0.06) node[below]{$-1/27$};
\draw (-0.06,{40/9})--(0.06,{40/9}) node[right]{$2/9$};
\draw (-0.06,{10/3})--(0.06,{10/3}) node[right]{$1/6$};
\fill ({-40/9},{40/9}) circle(1.6pt);
\fill (0,{10/3}) circle(1.3pt);
\node[above left,align=right] at ({-40/9},{40/9}) {three equal\\roots};
\node at (-1.55,2.25) {$\mathcal C_B$};
\node[below right] at (0,0) {$0$};
\end{tikzpicture}

\caption{The secondary corner of the principal-scale saturation domain, in coordinates \((\alpha,Bd)\).
The shaded interior corresponds to three distinct positive roots of
\(z^3-z^2+(3Bd/2)z+\alpha\). The solid boundary has a repeated root;
the dashed segment has a zero root. The dot at
\((-1/27,2/9)\) has three equal roots. The solid curve is
\((\alpha,Bd)=(-t^2(1-2t),2(2t-3t^2)/3)\), \(0\le t\le1/2\).
}
\label{fig:corner-limit-set}
\end{figure}

\begin{proof}[Proof of Proposition~\ref{corner:complete-cluster-set}]
Boundedness of $\alpha,d$ follows from
Proposition~\ref{corner:secondary-bounds}, so consider any cluster point.
For $d_0>0$, Lemma~\ref{corner:limiting-critical-states} gives
a cubic with positive critical states $z_-\leq z_+$,
nonnegative maximum, nonpositive minimum, and
$F_{\alpha_0,d_0}(0)=\alpha_0\leq0$.  Hence it has three
nonnegative real roots, counted with multiplicity.  This conclusion
includes equality of critical values or $\alpha_0=0$ by continuity,
or directly by the intermediate value theorem on the intervals
$[0,z_-]$, $[z_-,z_+]$, and $[z_+,\infty)$.
At $T_{3,0}=1/3$ the two critical states coincide; their two
value inequalities force a triple root there.
For $d_0=0$ the limiting polynomial is $z^2(z-1)$.
In every case Vieta's identities give membership in $\mathcal C_B$.

The discriminant of $z^3-z^2+T_3 z+\alpha$ is the expression
in \eqref{corner:discriminant}.  Conversely, that discriminant and
$\alpha\leq0$, $T_3\geq0$ give three real roots with sum one,
nonnegative product, and nonnegative pair sum.  If a root were
negative, the product condition would allow us to write the two
nonpositive roots as $-u,-v$, with $u>0$, $v\geq0$, and the third
as $1+u+v$.  Their pair sum would be
$-u^2-uv-v^2-u-v<0$, a contradiction.  This proves
\eqref{corner:discriminant}.  Strict positivity and distinctness
characterize the interior; a zero root or a repeated root is on its
boundary.

For nonnegative $r_i$ with sum one,
$r_1r_2r_3\leq1/27$ and
$r_1r_2+r_1r_3+r_2r_3\leq1/3$.
These inequalities prove \eqref{corner:sharp-cluster-bounds}.
Equality in the second requires $r_1=r_2=r_3=1/3$ and gives the
stated top endpoint.  If $a\geq0$ eventually, every cluster point
has $\alpha_0=0$.  Its polynomial is
$z(z^2-z+T_{3,0})$, whose two remaining roots are real and
nonnegative only if $T_{3,0}\leq1/4$.  This proves
\eqref{corner:nonnegative-offset-bound}.

For the interior converse, fix a compact
$\mathcal K\Subset\operatorname{int}\mathcal C_B$. The ordered positive
roots $r_j$ have uniform positive lower bounds, pairwise
separation, and nonzero derivative margins.  Lemma
\ref{corner:joint-quotient} transfers them to three simple roots on
disjoint positive $z$-intervals for every independently small positive
$\mu,\sigma$, uniformly on $\mathcal K$. The other three roots follow from the
charts of Lemma~\ref{corner:realization-charts}.  The first chart has its
negative root at $x=\alpha$, the second at $X=1$, and the third at
$X=1/C^2$, all with strict uniform derivative margins.
The negative secondary root and the three positive secondary roots
have size $\mu^2\sigma^2$, whereas the outer negative root has size
$\mu^2\sigma$ and the outer positive root size $\mu^2$. The six isolating
intervals are therefore disjoint for small $\sigma$.
All their states belong to $I$ for small scales.

The standard unfolding realizes \eqref{corner:weights} on one product
$(0,\mu_0)\times(0,\sigma_0)$ for $\mathcal K$,
and~\cite[Proposition~\ref*{main-prop:common-count}]{main2026} makes the six
simple roots exhaustive; the same charts give \eqref{corner:interior-roots}.

For any boundary point, convexly perturb its ordered root triple toward
$(1/6,1/3,1/2)$ to obtain interior parameters $(\alpha_j,d_j)$ tending
to it. For the $j$th interior parameter choose
positive $\mu_j,\sigma_j<1/j$ below the two existence thresholds
just proved.  The corresponding return map has six simple roots,
and its normalized parameters are exactly $(\alpha_j,d_j)$.
This proves the diagonal assertion and, together with necessity,
identifies the full cluster set as $\mathcal C_B$.
It also proves sharpness of the unrestricted cluster bounds.

\end{proof}

\subsection{Hausdorff convergence}

To obtain convergence of the full domains, we first place the normalized
parameters of all six-cycle returns in one common rectangle.
For the standard unfolding and the $\mathcal V,I$ of
Definition~\ref{def:actual-secondary-domain},
Proposition~\ref{corner:secondary-bounds} bounds the full local domain by
$-K\le\alpha\le K\mu^2$ and $0<d\le K$. Choose
$\mathcal B=[-M,M]^2$, $M>K$, containing $\mathcal C_B$ in its interior.
The weighted parameters~\eqref{corner:weights} tend uniformly
to zero on this rectangle for independently small positive scales. Thus
$\Lambda\circ w_{\mu,\sigma}$ is defined throughout $\mathcal B$ after
one common reduction, and
$\mathscr S_{\mu,\sigma}\subset\operatorname{int}\mathcal B$.
The rectangle therefore contains the normalized parameters of every six-cycle
return in $\mathcal V$, including those with positive offset.

\begin{proof}[Completion of Proposition~\ref{thm:saturation-geometry}]
We apply the compactness and finite-net argument that completes the proof
of~\cite[Theorem~\ref*{main-thm:main-saturation-summary}]{main2026} in
\cite[Section~\ref*{main-subsec:four-root-limit}]{main2026}. The common rectangle above makes
every closure compact. An interior singleton of $\mathcal C_B$ is realized
by part~2 of
Proposition~\ref{corner:complete-cluster-set} for every sufficiently small
independent scale pair, so these closures are nonempty.

For the upper inclusion, the weighted parameters tend uniformly to the fixed
reference on $\mathcal B$, so part~1 applies along arbitrary sequences
$\mu_n,\sigma_n\to0^+$. For the lower inclusion, the root perturbation
above gives interior density. Every finite interior net is compact, so
part~2 realizes it in the standard unfolding on one product of positive
scale intervals, with no relation between the scales. These are the
two-scale inputs to the cited argument, which gives both Hausdorff inclusions.
\end{proof}
